\section{Kiểm thử và đánh giá hệ thống}

\subsection{Chiến lược và Quy trình kiểm thử}
Nhằm đảm bảo chất lượng, độ tin cậy và tính toàn vẹn của nền tảng UniConnect trước khi đưa vào vận hành thực tế, đồ án thiết lập quy trình kiểm thử phần mềm chuẩn hóa tuân thủ theo tiêu chuẩn IEEE 829\cite{ieee829test}. Chiến lược kiểm thử được xây dựng theo mô hình Kim tự tháp kiểm thử (Test Pyramid), bao gồm nhiều cấp độ kiểm tra chặt chẽ:
\begin{itemize}
    \item \textbf{Kiểm thử đơn vị và tích hợp Backend (Unit \& Integration Tests):} Xác minh tính đúng đắn của từng hàm nghiệp vụ, cấu hình bảo mật, giao dịch cơ sở dữ liệu và các ràng buộc bất biến (invariants) giữa các phân hệ.
    \item \textbf{Kiểm thử vòng đời kết nối và phục hồi (Lifecycle \& Resilience Tests):} Kiểm tra chu kỳ sống của phiên làm việc cơ sở dữ liệu, tính toàn vẹn khi rollback và khả năng tự phục hồi kết nối.
    \item \textbf{Kiểm thử giao diện người dùng (Frontend Component Tests):} Đảm bảo các thành phần giao diện hiển thị chính xác, xử lý trạng thái và lưu trữ token an toàn.
    \item \textbf{Đánh giá tối ưu hóa hiệu năng (Performance Optimization Evaluation):} Đo lường hiệu quả của kỹ thuật phân tách mã nguồn (Code-splitting) trên kích thước bundle truyền tải qua mạng.
    \item \textbf{Kiểm thử chấp nhận người dùng (User Acceptance Testing - UAT):} Thực nghiệm các kịch bản người dùng thực tế trên môi trường container tương đương production.
\end{itemize}

\subsection{Kiểm thử tự động Backend (Automated Backend Testing)}
Hệ thống backend của UniConnect được kiểm thử tự động toàn diện bằng framework \texttt{pytest} kết hợp thư viện \texttt{pytest-asyncio} và client HTTP bất đồng bộ \texttt{httpx.AsyncClient}. Bộ kiểm thử sử dụng cơ chế session-scoped database fixture với SQLite in-memory và transaction rollback độc lập cho từng ca kiểm thử, loại bỏ hoàn toàn sự phụ thuộc chéo giữa các bài test.

Toàn bộ \textbf{64 ca kiểm thử tự động} được phân bổ trong 11 bộ kiểm thử (Test Suites) chuyên biệt tại thư mục \texttt{server/tests/}, đạt tỷ lệ vượt qua 100\% (64/64 passed). Bảng~\ref{tab:backend_test_suites} tổng hợp chi tiết kết quả kiểm thử của từng phân hệ:

\begingroup
\singlespacing
\small
\renewcommand{\arraystretch}{1.25}
\begin{xltabular}{\textwidth}{|>{\raggedright\arraybackslash}p{2.8cm}|>{\raggedright\arraybackslash\ttfamily\small}p{3.8cm}|>{\centering\arraybackslash}p{1.1cm}|>{\raggedright\arraybackslash}X|}
\caption{Bảng tổng hợp kết quả 64 ca kiểm thử tự động backend UniConnect} \label{tab:backend_test_suites} \\
\hline
\textbf{Bộ kiểm thử} & \textbf{Tập tin mã nguồn} & \textbf{Số test} & \textbf{Nghiệp vụ \& Invariant được kiểm tra} \\ \hline
\endfirsthead

\multicolumn{4}{c}{{\bfseries Bảng \thetable{} -- Tiếp theo từ trang trước}} \\
\hline
\textbf{Bộ kiểm thử} & \textbf{Tập tin mã nguồn} & \textbf{Số test} & \textbf{Nghiệp vụ \& Invariant được kiểm tra} \\ \hline
\endhead

\hline
\multicolumn{4}{r}{\textit{Còn tiếp ở trang sau...}} \\
\endfoot

\hline
\endlastfoot

Quản trị \& Kiểm duyệt & \texttt{test\_\allowbreak admin\_\allowbreak command\_\allowbreak center.py} & 10 & Dashboard KPI tổng hợp, kiểm duyệt \& ẩn hoạt động vi phạm, quy trình xét duyệt hồ sơ tổ chức giáo dục, phân quyền RBAC đa cấp, bảo vệ trạng thái dừng (terminal states), xuất dữ liệu sinh viên CSV UTF-8 BOM. \\ \hline
Sức khỏe API \& Giám sát & \texttt{test\_\allowbreak api.py} & 3 & Endpoint \texttt{/health} giám sát độ trễ CSDL, gắn HTTP timing header, cô lập an toàn lỗi 503 khi CSDL quá tải mà không rò rỉ thông tin nhạy cảm. \\ \hline
Điểm danh \& Chứng nhận CTXH & \texttt{test\_\allowbreak attendance\_\allowbreak trophy.py} & 6 & Khoảng cách Haversine, sinh \& thẩm định token QR xoay động 30s (HMAC-SHA256), chuẩn hóa mã chứng nhận CTXH, chống điểm danh lặp (idempotency), từ chối GPS có sai số $acc > 50$m. \\ \hline
Lịch thông minh \& Xung đột & \texttt{test\_\allowbreak calendar.py} & 9 & Đối soát xung đột lịch nửa mở $[start, end)$, bận định kỳ hàng tuần, ngoại lệ ngày nghỉ (vacation), xung đột mềm khi đăng ký tham gia, tính năng xem trước dời lịch (reschedule preview). \\ \hline
Trợ lý AI \& Tìm kiếm & \texttt{test\_\allowbreak chat\_\allowbreak assistant.py} & 5 & Function calling tìm kiếm hoạt động và nhóm, bảo đảm quyền riêng tư lịch bận cá nhân, cơ chế fallback khi Gemini API ngoại tuyến, kiểm soát tần suất hội thoại (Rate Limiting). \\ \hline
Chu kỳ sống CSDL \& Rate Limiter & \texttt{test\_\allowbreak database\_\allowbreak observability.py} & 7 & Tham số Connection Pool (\texttt{pool\_size=20}, \texttt{max\_overflow=10}, \texttt{pool\_recycle=1800}), commit/rollback tự động của \texttt{get\_db()}, Rate Limiter dọn dẹp sliding-window và giới hạn bộ nhớ. \\ \hline
Luồng tích hợp đầu cuối (E2E) & \texttt{test\_\allowbreak e2e\_\allowbreak critical\_\allowbreak flows.py} & 4 & Vòng đời trọn vẹn: Tạo sự kiện $\rightarrow$ Điểm danh $\rightarrow$ Ghi nhận CTXH $\rightarrow$ Danh hiệu; Độc lập giữa xung đột lịch và đăng ký; Ủy nhiệm quyền Co-host nhóm; Xét duyệt vai trò tổ chức kèm audit log. \\ \hline
Quản lý Nhóm \& Đồng tổ chức & \texttt{test\_\allowbreak groups\_\allowbreak cohosting.py} & 9 & Phân quyền RBAC nhóm sinh viên, phê duyệt yêu cầu gia nhập, phân tách quyền ủy nhiệm \texttt{can\_open\_checkin} và \texttt{can\_manage\_activity}, giới hạn số nhóm co-host tối đa, rollback giao dịch khi lỗi. \\ \hline
Tương tác \& Thông báo & \texttt{test\_\allowbreak interactions\_\allowbreak notifications.py} & 7 & Bình luận phân cấp đa tầng (\texttt{parent\_id}), tính toàn vẹn Like duy nhất, tương thích ngược Pydantic Schemas, phát sinh thông báo tương tác kèm liên kết điều hướng. \\ \hline
Khởi tạo Ứng dụng & \texttt{test\_\allowbreak main.py} & 1 & Khởi động thành công ứng dụng FastAPI, nạp middleware CORS, định tuyến router và cấu hình lifecycle handlers. \\ \hline
Hồ sơ \& Thống kê Sinh viên & \texttt{test\_\allowbreak profile\_\allowbreak stats.py} & 3 & Đối soát ngưỡng điểm CTXH, tính toán danh hiệu/thứ hạng (rank/title), cơ chế bảo vệ quyền riêng tư thông tin hồ sơ người dùng công khai. \\ \hline
\multicolumn{2}{|l|}{\textbf{Tổng số ca kiểm thử}} & \textbf{64} & \textbf{Tỷ lệ vượt qua: 64/64 (100\% Passed)} \\ \hline
\end{xltabular}
\endgroup

\subsubsection{Đặc tả kỹ thuật về Kiểm thử Chu kỳ sống Connection Pool}
Trong phân hệ quản trị kết nối dữ liệu (\texttt{test\_database\_observability.py}), hệ thống thực hiện kiểm thử cấu hình và lifecycle của connection pool với các tham số cấu hình production:
\begin{itemize}
    \item Tham số thiết lập: \texttt{pool\_size = 20}, \texttt{max\_overflow = 10}, \texttt{pool\_recycle = 1800}, \texttt{pool\_pre\_ping = True}.
    \item Kiểm thử xác nhận: Bộ kiểm thử mô phỏng chu kỳ mở kết nối, cấp phát phiên (session checkout), thực thi truy vấn nghiệp vụ, giải phóng kết nối về hàng đợi (session return), và kiểm tra tính toàn vẹn khi xảy ra sự kiện hủy kết nối đột ngột (DB restart resilience).
    \item \textbf{Lưu ý phương pháp luận:} Kết quả kiểm thử này chứng minh tính đúng đắn về mặt cấu hình, tính an toàn của chu kỳ sống kết nối và khả năng tự phục hồi (Lifecycle and Resilience Test). Đây là kiểm thử hồi quy logic kiến trúc, không phải là thử nghiệm đánh giá hiệu năng chịu tải dưới áp lực đồng thời cực hạn (Load/Concurrency Benchmark).
\end{itemize}

\subsection{Kiểm thử tự động Giao diện Frontend (Frontend Component Testing)}
Giao diện người dùng Web Client được kiểm thử tự động thông qua công cụ Vitest kết hợp cùng thư viện React Testing Library và môi trường mô phỏng JSDOM. Các ca kiểm thử tập trung vào tính tương tác và tính nhất quán dữ liệu của client:
\begin{itemize}
    \item \textbf{Kiểm thử xác thực và quản lý token:} Kiểm tra tính năng lưu trữ Access Token trong bộ nhớ mã nguồn và làm mới phiên làm việc qua HTTP-only Refresh Token.
    \item \textbf{Kiểm thử hiển thị bản đồ tương tác:} Kiểm tra việc dựng bản đồ khuôn viên Leaflet, gắn nhãn các điểm ghim hoạt động dựa trên tọa độ địa lý và xử lý sự kiện click mở thông tin chi tiết.
    \item \textbf{Kiểm thử bộ điều hướng và bảo vệ tuyến đường (Route Guards):} Đảm bảo các trang chức năng yêu cầu quyền (như trang Quản trị viên, trang Tổ chức sự kiện) chặn điều hướng và chuyển hướng chính xác về màn hình đăng nhập khi người dùng chưa xác thực.
\end{itemize}

\subsection{Đánh giá Tối ưu hóa Hiệu năng Frontend (Bundle Optimization)}
Trong môi trường mạng học đường di động (3G/4G), kích thước gói tài nguyên JavaScript ban đầu ảnh hưởng quyết định đến thời gian hiển thị nội dung đầu tiên (First Contentful Paint - FCP). Nhằm nâng cao trải nghiệm sinh viên, hệ thống đã ứng dụng kỹ thuật phân tách mã nguồn động (Code-splitting) thông qua cơ chế \texttt{React.lazy()} và cấu hình chia nhỏ gói (Chunking Strategy) của Vite/Rollup. 

Bảng~\ref{tab:bundle_optimization} thể hiện dữ liệu đo lường thực tế kích thước gói mã nguồn trước và sau khi thực hiện tối ưu hóa:

\begin{table}[H]
\centering
\small
\begin{tabular}{|l|c|c|c|}
\hline
\textbf{Chỉ số đo lường} & \textbf{Trước tối ưu (Single Bundle)} & \textbf{Sau tối ưu (Code-splitting)} & \textbf{Mức độ tối ưu} \\ \hline
Kích thước Entry JS thô (Raw Size) & 842.80 kB & 236.02 kB & \textbf{Giảm 72.0\%} \\ \hline
Kích thước Entry JS nén (Gzip Size) & 240.26 kB & 72.15 kB & \textbf{Giảm 70.0\%} \\ \hline
Số lượng phân mảnh (Dynamic Chunks) & 1 tập tin nguyên khối & 8 phân mảnh theo tuyến & Tải theo nhu cầu \\ \hline
Thời gian hoàn tất tải trang (3G mạng chậm) & $\approx$ 4.2 giây & $\approx$ 1.3 giây & \textbf{Nhanh hơn 3.2 lần} \\ \hline
\end{tabular}
\caption{Kết quả đo lường tối ưu hóa kích thước gói mã nguồn Frontend}
\label{tab:bundle_optimization}
\end{table}

Kết quả đo lường khẳng định việc giảm tới 72\% kích thước file mã nguồn ban đầu giúp ứng dụng khởi động tức thì, tối ưu hóa mức tiêu thụ băng thông và mang lại trải nghiệm mượt mà cho người dùng ngay cả trong điều kiện mạng yếu tại các khu vực xa trung tâm giảng đường.

\subsection{Ma trận Kiểm thử Chấp nhận Người dùng (User Acceptance Testing - UAT)}
Nhằm thiết lập bộ tiêu chuẩn đối soát nghiệm thu các yêu cầu nghiệp vụ của hệ thống, đồ án xây dựng Ma trận kiểm thử chấp nhận người dùng (User Acceptance Testing Matrix) gồm 10 kịch bản cốt lõi. Các kịch bản này bao quát toàn bộ các nhóm vai trò (Người dùng, Ban điều hành Nhóm, Quản trị viên và Khách vãng lai), định hình các điều kiện tiên quyết, hành động thực hiện và tiêu chuẩn nghiệm thu tương ứng. Chi tiết được trình bày trong Bảng~\ref{tab:uat_matrix}.

\begingroup
\singlespacing
\small
\renewcommand{\arraystretch}{1.2}
\begin{xltabular}{\textwidth}{|>{\centering\arraybackslash}p{1.4cm}|>{\raggedright\arraybackslash}p{1.8cm}|>{\raggedright\arraybackslash}X|>{\raggedright\arraybackslash}X|>{\centering\arraybackslash}p{1.6cm}|}
\caption{Ma trận kiểm thử chấp nhận người dùng (User Acceptance Testing Matrix)} \label{tab:uat_matrix} \\
\hline
\textbf{Mã kịch bản} & \textbf{Tác nhân} & \textbf{Hành động thực hiện} & \textbf{Kết quả kỳ vọng} & \textbf{Tiêu chí nghiệm thu} \\ \hline
\endfirsthead

\multicolumn{5}{c}{{\bfseries Bảng \thetable{} -- Tiếp theo từ trang trước}} \\
\hline
\textbf{Mã kịch bản} & \textbf{Tác nhân} & \textbf{Hành động thực hiện} & \textbf{Kết quả kỳ vọng} & \textbf{Tiêu chí nghiệm thu} \\ \hline
\endhead

\hline
\multicolumn{5}{r}{\textit{Còn tiếp ở trang sau...}} \\
\endfoot

\hline
\endlastfoot

UAT-01 & Người dùng & Nhập email, mật khẩu hợp lệ; nhấn Đăng nhập. & Cấp JWT token, duy trì phiên làm việc an toàn, chuyển hướng về Bản đồ. & Đạt chuẩn \\ \hline
UAT-02 & Người dùng & Mở màn hình chính; bật định vị GPS; chọn bộ lọc thể thao. & Bản đồ hiển thị các điểm ghim hoạt động thể thao trong bán kính khuôn viên trường. & Đạt chuẩn \\ \hline
UAT-03 & Người dùng & Nhấn "Tham gia" sự kiện có thời gian trùng với lịch cá nhân. & Hệ thống phát hiện giao nhau thời gian và hiển thị cảnh báo trùng lịch (\texttt{hard\_conflict}). & Đạt chuẩn \\ \hline
UAT-04 & Người dùng & Quét mã QR xoay động tại bàn check-in sự kiện trong chu kỳ 30s. & Backend xác thực chữ ký HMAC-SHA256 hợp lệ, ghi nhận tham gia thành công. & Đạt chuẩn \\ \hline
UAT-05 & Người dùng & Đứng tại hội trường ($d \le 100$m, $acc \le 50$m); nhấn "Điểm danh GPS". & Hệ thống tính khoảng cách không gian thỏa mãn; cập nhật \texttt{attendance\_confirmed=True}. & Đạt chuẩn \\ \hline
UAT-06 & Người dùng & Thử điểm danh GPS khi ở ngoài khuôn viên ($d > 100$m). & Hệ thống từ chối ghi nhận và phản hồi thông báo vị trí nằm ngoài phạm vi cho phép. & Đạt chuẩn \\ \hline
UAT-07 & Người dùng & Hoàn tất điểm danh sự kiện tình nguyện có định mức 0.5 ngày CTXH. & Tự động tích lũy thêm 0.5 ngày CTXH; mở khóa huy hiệu Trophy tương ứng. & Đạt chuẩn \\ \hline
UAT-08 & Ban tổ chức & Nhóm chủ trì gửi lời mời Nhóm khác cùng Đồng tổ chức (Co-hosting). & Nhóm đối tác nhận lời mời ở trạng thái \texttt{PENDING} và có thể chấp thuận/từ chối. & Đạt chuẩn \\ \hline
UAT-09 & Co-host & Thành viên Nhóm đối tác (đã chấp thuận Co-host) mở phiên điểm danh sự kiện. & Động cơ phân quyền RBAC xác nhận quyền ủy nhiệm \texttt{can\_open\_checkin} thành công. & Đạt chuẩn \\ \hline
UAT-10 & Quản trị viên & Admin xử lý báo cáo vi phạm nội dung và nhấn "Xuất danh sách người dùng". & Gỡ bỏ nội dung vi phạm tức thì; tải về tập tin CSV định dạng chuẩn UTF-8 BOM. & Đạt chuẩn \\ \hline
\end{xltabular}
\endgroup
